\documentclass[preview, border=12pt]{standalone}
\input{/Users/gcrescenzo/Documents/School/LaTeX Documents/latex-class-notes/preamble-hw.tex}
\input{/Users/gcrescenzo/Documents/School/LaTeX Documents/latex-class-notes/makros.tex}
\renewcommand*{\mathbf}[1]{\varmathbb{#1}}
\renewcommand*{\emph}[1]{\textit{#1}}
\setlength{\parskip}{0.5\baselineskip}
\begin{document}
Let $X = \{1,\ldots ,n\}$. For each $\alpha \in \{1,\ldots ,n\}$, let $\hat{\alpha}:C(X) \rightarrow \bfC$ be given by $\hat{\alpha}(f) := f(\alpha)$. Then $\Sigma_{C(X)} = \{\hat{1},\hat{2},\ldots ,\hat{n}\}$, where $\Sigma_{C(X)}$ is the set of all nonzero algebra homomorphisms from $C(X)$ to $\bfC$ (i.e., the characters of $C(X)$). \\

First, note that $\{\delta_1,\ldots ,\delta_n\}$ is a basis for $C(X)$, where each $\delta_i$ is defined by:
\[
	\delta_i(j) = \begin{cases} 1, & i = j\\ 0, & \text{otherwise}\end{cases}.
\] 
Clearly $\{\hat{1},\hat{2},\ldots ,\hat{n}\} \subseteq \Sigma_{C(X)}$, as for any $\hat{\alpha} \in \{\hat{1},\hat{2},\ldots ,\hat{n}\}$ we have:
\begin{align*}
	\hat{\alpha}(f+g) &= (f+g)(\alpha) = f(\alpha) + g(\alpha), \\
	\hat{\alpha}(fg) &= (fg)(\alpha) = f(\alpha) g(\alpha), \\
	\hat{\alpha}(\lambda f) &= (\lambda f)(\alpha) = \lambda f(\alpha).
\end{align*}
Moreover, each $\hat{\alpha}$ is nonzero, as $\hat{\alpha}(\delta_\alpha) = \delta_\alpha(\alpha) = 1$. \\

Claim: for any $h \in \Sigma_{C(X)}$, there is exactly one $\delta_i \in \{\delta_1,\ldots ,\delta_n\}$ such that $h(\delta_i) = 1$, and $h(\delta_j) = 0$ for every $j \neq i$. Since $\delta_i^2 = \delta_i$, we have $h(\delta_i)^2 = h(\delta_i^2) = h(\delta_i)$. Hence $h(\delta_i) = 1$ or $0$. If it were the case that $h(\delta_i) = 0$ for all $i = 1,\ldots ,n$, then $h$ would be the zero functional which contradicts it being an element of $\Sigma_{C(X)}$. Now consider $\sum_{\alpha = 1}^{n} \delta_\alpha$. It is also true that $\left(\sum_{\alpha = 1}^{n} \delta_\alpha\right)^2 = \sum_{\alpha =1}^{n} \delta_\alpha$, hence:
\[
	\left( \sum_{\alpha = 1}^{n} h(\delta_\alpha) \right)^2 = \sum_{\alpha = 1}^{n} h(\delta_\alpha).
\]
It must be the case that $\sum_{\alpha = 1}^{n} h(\delta_\alpha) = 1$. Since $h(\delta_\alpha)$ can be either $1$ or $0$, this is only true when $h(\delta_\alpha) = 1$ for exactly one $\delta_\alpha \in \{\delta_1,\ldots ,\delta_n\}$. \\

With this, let $h \in \Sigma_{C(X)}$. Suppose $i$ is the unique index such that $h(\delta_i) = 1$. Then for any $f \in C(X)$:
\begin{align*}
	h(f)
	&= h \left( \sum_{\alpha = 1}^{n}f(\alpha)\delta_\alpha  \right)  \\
	&= \sum_{\alpha = 1}^{n} f(\alpha)h(\delta_\alpha) \\
	&= f(i) \\
	&= \hat{i}(f).
\end{align*}
Thus $h \in \{\hat{1},\hat{2},\ldots ,\hat{n}\}$, establishing $\Sigma_{C(X)} = \{\hat{1},\hat{2},\ldots ,\hat{n}\}$.
\end{document}
